\begin{beginnerstatement}[Kurz erklärt: Qualitätsziele]

Testbarkeit bedeutet, dass sich ein Baustein mit klaren Prüfungen kontrollieren
lässt, möglichst unabhängig von anderen Teilen. Wartbarkeit beschreibt, wie
gut Software verstanden, geändert und aktualisiert werden kann.
Reproduzierbarkeit bedeutet, dass gleiche festgelegte Eingaben und
Einstellungen zu nachvollziehbaren Ergebnissen führen. Im Template sind dies
Qualitätsziele, an denen die spätere technische Prüfung ausgerichtet wird.

LOC bedeutet \emph{Lines of Code}, also Codezeilen. Ein Quality Gate ist eine
automatische Schranke: Es sammelt Prüfergebnisse und liefert einen Fehlercode,
wenn eine blockierende Regel verletzt ist. Eine Severity beschreibt die
Schwere eines Befunds. Warnungen weisen hier auf Risiken hin, ohne allein zu
blockieren; ERROR blockiert regulär. Die Zahlen sind Projektregeln und keine
Garantie dafür, dass kleiner Code automatisch leicht verständlich oder gut
wartbar ist.

\end{beginnerstatement}
